\section{A portable measured-coordinate experiment}
The fixture \texttt{psm\_temperature\_2500} contains 680 operating points,
340 each at rotor-temperature references 30 and 70 $^\circ$C, at a requested
2500 rpm. The stored machine code is PSM and the phase selector is three-phase.
The nominal 1500-rpm variant changes six speed-related signals while retaining
all current and torque values; it is excluded as an independent experiment.

Each operating point groups logger records by requested speed, excitation,
temperature reference and measurement counter. Measured mean currents and
voltages, medians, sample deviations and counts are retained. Flux is the
stationary inversion of group means using recorded mean electrical speed and
the recorded resistance. No residual-based deletion or extra settling window
is introduced. The acquisition timing does not justify treating three stored
records as three independent calibrated repetitions.

Within each temperature slice, construct a vector-valued linear Delaunay
interpolant on measured currents. For positive q representatives above
10\% of the slice maximum $|i_q|$, evaluate its mirror at $(i_d,-i_q)$.
Keep only mirrors inside the convex hull. This addresses numerical current
matching explicitly, but introduces correlated interpolation error.
Projection is applied to the two evaluated values; we do not claim that a new
interpolant through projected vertices would automatically retain exact parity.
The raw field and removed component stay available.
